\section{Discussion}

\subsection{Hypothesis Evaluation}

Our pre-registered hypothesis stated that a batch-64 state transition
circuit at depth~32 could be generated in under 60~seconds with fewer
than 100,000~R1CS constraints. The results require nuanced evaluation:

\textbf{Constraint target (100K): Rejected.} Batch~64 at depth~32
requires 2,192,256~constraints---21.9$\times$ over the 100K target.
Even with the most aggressive optimization (circomlib's
\texttt{SMTProcessor} at approximately 9,000 constraints per transaction),
the minimum is approximately 576,000~constraints. The 100K target was
based on an incorrect assumption about the relationship between constraint
count and proving feasibility; in practice, modern Groth16 provers
handle millions of constraints within acceptable time budgets.

\textbf{Proving time target (60s): Partially confirmed.}
\begin{itemize}
\item SnarkJS (JavaScript, commodity hardware): estimated 120--143~seconds.
  \textbf{Rejected} (2.0--2.4$\times$ over target).
\item Rapidsnark (C++, commodity hardware): estimated 14--35~seconds.
  \textbf{Confirmed} with 1.7--4.3$\times$ margin.
\item Rapidsnark (C++, server hardware): estimated 8--15~seconds.
  \textbf{Confirmed} with wide margin.
\item GPU-accelerated (ICICLE-Snark): estimated 3--8~seconds.
  \textbf{Confirmed}.
\end{itemize}

The overall verdict is that the engineering goal---batch-64 state
transition proofs in under 60~seconds---is achievable, but \emph{only}
with production-grade native provers. SnarkJS is suitable for development
and testing, not production deployment.

\subsection{Constraint Formula Significance}

The exact formula $C_{\text{tx}} = 1{,}038 \cdot (d+1)$ has direct
engineering implications:

\begin{enumerate}
\item \textbf{Capacity planning}: Given a proving time budget $T$ and a
  prover's throughput in constraints per second, the maximum batch size
  is $n_{\max} = \lfloor T \cdot \text{throughput} / C_{\text{tx}} \rfloor$.
  For Rapidsnark at 400K constraints/second and $T = 30$s:
  $n_{\max} = \lfloor 12{,}000{,}000 / 34{,}254 \rfloor = 350$ transactions
  at depth~32.

\item \textbf{Depth selection}: Reducing depth from 32 to 20 decreases
  per-transaction cost by $1 - (21/33) = 36\%$, enabling proportionally
  larger batches. Depth~20 still supports $2^{20} \approx 1$M unique
  keys, sufficient for most enterprise deployments.

\item \textbf{Circuit predictability}: The zero-residual fit means that
  the Circom compiler introduces no unexpected constraint overhead
  for this circuit structure. Constraint counts can be predicted
  exactly before compilation.
\end{enumerate}

\subsection{Comparison with Production Systems}

Our per-transaction constraint count of 34,254 at depth~32 is
0.59$\times$ the Hermez rollup figure of approximately 58,000 per
transaction~\cite{hermez2021circuits}. This is consistent: Hermez
includes EdDSA signature verification ($\sim$6,000 constraints),
token balance updates, fee processing, and nonce management that
our circuit omits. The delta of $\sim$24,000 constraints represents
the cost of full transaction validation beyond Merkle state updates.

Relative to Tornado Cash~\cite{tornadocash2019} (5,000--7,000
constraints for a single depth-20 Merkle proof), our per-transaction
cost at depth~20 (21,798 constraints) is 3.1--4.4$\times$ higher,
consistent with the factor of performing \emph{two} Merkle path
computations (old and new) plus leaf hashing at each of the 21
effective levels.

\subsection{Optimization Paths}

Several optimization strategies could reduce the constraint count:

\textbf{SMTProcessor integration.} Circomlib's \texttt{SMTProcessor}
shares intermediate computations between the old-root verification and
new-root computation, potentially reducing per-transaction cost by
up to 40\%. However, it introduces state machine complexity (insert,
update, delete modes) that is unnecessary for an update-only circuit.
For the enterprise validium use case, the simplicity and auditability
of the dual-path approach may outweigh the constraint savings.

\textbf{Batch-level sibling sharing.} When multiple transactions in a
batch update keys that share common Merkle path prefixes, the
corresponding sibling hashes are identical. A circuit-level optimization
could detect and reuse these shared siblings. In the worst case (random
keys), no sharing occurs; in enterprise scenarios with sequential keys,
significant sharing is possible. This optimization complicates the circuit
but could reduce effective constraints by 20--50\% for favorable
access patterns.

\textbf{Reduced tree depth.} Depth~20 provides $2^{20} \approx 1$M
key slots at 36\% fewer constraints. For MVP deployments where the
key space is bounded, this is the simplest optimization.

\textbf{Proof aggregation.} Rather than increasing batch size, multiple
smaller proofs can be aggregated using recursive SNARKs. This approach
decouples per-batch proving time from total throughput: four parallel
batch-16 proofs aggregated into a single recursive proof could achieve
batch-64 effective throughput with each individual proof requiring only
548K constraints.

\subsection{Limitations}

\textbf{Single-run measurements.} Each configuration was benchmarked
once. While the constraint counts are deterministic (verified by the
exact formula fit), proving times exhibit variance due to OS scheduling,
memory allocation, and Powers of Tau ceremony size. A full stochastic
analysis with 30+ repetitions per configuration would strengthen the
proving time estimates.

\textbf{SnarkJS only.} All proving time measurements use SnarkJS, a
pure JavaScript prover. Production proving time estimates for Rapidsnark
and GPU-accelerated provers are extrapolated from published benchmarks,
not measured directly. The 4--10$\times$ speedup range is broad.

\textbf{Update-only semantics.} The circuit assumes all transactions
are updates to existing key-value pairs. Insertions (new keys) and
deletions require different Merkle proof structures not covered by
this circuit.

\textbf{No signature verification.} Unlike production rollup circuits
(Hermez, Polygon zkEVM), our circuit does not verify transaction
authorization. In the enterprise validium model, authorization is
handled by the permissioned node, not the ZK circuit.
